% poc/g3-protrusion.tex — G3-bis: pdfTeX によるぶら下げ組（行末突出し）
% の可否判定 PoC。結論: 可能（条件つき）。根拠は各パネルの \typeout
% マーカと \rightmarginkern 読み戻し値（ログ証拠）＋同梱 PDF の行構造。
%
% 背景: 先行メモは「pdfTeX には行末突出し機構がないため不可能（確定制
% 限）」としていたが、pdfTeX には \pdfprotrudechars=2 と \rpcode/\lpcode
% による文字突出し（microtype の基盤）がある。本 PoC はバックエンドと
% 同一のフォント系（bxcjkjatype[ipaex-type1,noCJKpunct] + always-on）で
% 以下を検証する:
%   (a) 行末ちょうどに来た 、 が \rpcode で視覚的に突出するか
%   (b) バックエンドのノード列（ペアグルー/禁則ペナルティ/カーン/箱）が
%       突出しを阻害しないか
%
% 判明した機構（pdfTeX ソース pdftex.web で裏付け）:
%   - 突出しは「margin kern」が行末文字の側に挿入され、行内グルーの
%     再配分（両端揃えの伸長）によって実現する。パック時の行端走査
%     （cp_skipable）は、非ゼロ幅の「明示」カーン（\kern: subtype=exp-
%     licit）や空でない箱を越えて行端文字に到達できない。
%   - カーン直後のブレイク（kern_break）は「直後が glue」で成立し、
%     パック時にそのカーンは 0 化され走査上 skipable になる。
%     「クラス幅 < グリフ幅」の詰めに使える唯一の形。
%   - \rpcode の値は quad（1em）の千分率 ⇒ rpcode=1000 で margin kern
%     = -1em（全角グリフちょうど全幅分）。
%
% パネル構成（emulated ノード列。実装は jlreq-pdftex-rules.sty の
% \ifjlreqpdf@hang ブロック）:
%   A: 裸 、 + 補正カーン + グルー + CJK マーカ OFF + rpcode → 期待 YES
%   B: A + CJK マーカーカーン（-2sp/+2sp） → 期待 NO（マーカが行端走査を
%      止める）⇒ ぶら下げ時にマーカを無効化する根拠
%   C: 箱詰めグリフ（現行 F2 の幅補正方式）+ rpcode → 期待 NO（箱が走査を
%      止める）⇒ 既定バリアントでは突出ししない（設計意図どおり）
%   D: 補正カーンとグルーの間に whatsit（\pdfsavepos） → 期待 NO
%      （ブレイクが glue に移りカーンが残る）
%   E: 禁則ペナルティ 10000 前置 → 期待 YES（penalty は skipable）
%
% 実行: pdflatex -interaction=nonstopmode poc/g3-protrusion.tex
% 判定: ログの G3-PANEL-X=YES/NO（YES=右マージンカーン非ゼロ）が期待と
% 一致すること。PDF の該当行の 、 が行長を 0.5zw だけ超えて見える。
\documentclass{article}
\usepackage[ipaex-type1,noCJKpunct]{bxcjkjatype}
\makeatletter
\let\CJK@global\global
\let\CJK@active\global
\let\CJK@local\global
\CJK@envStart{}{UTF8}{ipxm}%
\CJKspace
\makeatletter
\pdfprotrudechars=2
\parindent=0pt

\newdimen\zw \zw=10pt

% \panel{<id>}{<body>}: hsize=95pt の段落を組み、行端（最終から2行目=
% 、 が来る側）の右マージンカーンを \rightmarginkern で読んで判定し、
% G3-PANEL-<id>=YES/NO (<value>) をログに出す。
\def\panel#1#2{%
  \setbox10=\vbox{\hsize=95pt #2\par}%
  % 行ボックスの取出しはコピーから行う（原本から \lastbox で抜くと
  % 組上がりが失われて出荷されなくなる）
  \setbox20=\vbox{\unvcopy10
    \global\setbox11=\lastbox \unskip\unpenalty
    \global\setbox12=\lastbox}%
  \edef\g@rmk{\rightmarginkern12}%
  \ifnum\pdfstrcmp{\g@rmk}{0pt}=0
    \typeout{G3-PANEL-#1=NO (rightmargin=\g@rmk)}%
  \else
    \typeout{G3-PANEL-#1=YES (rightmargin=\g@rmk)}%
  \fi
  \box10\penalty0\relax}

\begin{document}
% フォントと rpcode: 、(U+3001=plane30 slot1), 。(slot2)
\setbox0=\hbox{、。漢}
\expandafter\let\expandafter\pfont\csname C70/ipxm/m/n/10/30\endcsname
\typeout{G3-FONT=\fontname\pfont}
\rpcode\pfont1=1000 \rpcode\pfont2=1000

% ---- Panel A: 本実装の原理形（裸 char＋補正カーン＋グルー、マーカ OFF）
\panel{A}{%
  \def\CJKglue{}\def\CJK@kern{}\def\CJK@CJK{}%
  漢\hskip 0pt plus 3pt 漢\hskip 0pt plus 3pt 漢\hskip 0pt plus 3pt
  漢\hskip 0pt plus 3pt 漢\hskip 0pt plus 3pt 漢\hskip 0pt plus 3pt
  漢\hskip 0pt plus 3pt 漢\hskip 0pt plus 3pt 漢
  、\kern -0.5\zw \hskip 0.5\zw plus 0.25\zw minus 0.5\zw
  漢漢漢漢漢漢漢漢漢}

% ---- Panel B: CJK マーカーカーンを残したまま
\panel{B}{%
  \def\CJKglue{}%
  漢\hskip 0pt plus 3pt 漢\hskip 0pt plus 3pt 漢\hskip 0pt plus 3pt
  漢\hskip 0pt plus 3pt 漢\hskip 0pt plus 3pt 漢\hskip 0pt plus 3pt
  漢\hskip 0pt plus 3pt 漢\hskip 0pt plus 3pt 漢
  、\CJK@kern \kern -0.5\zw \hskip 0.5\zw plus 0.25\zw minus 0.5\zw
  漢漢漢漢漢漢漢漢漢}

% ---- Panel C: 箱詰めグリフ（F2 現行方式）＋ rpcode
\panel{C}{%
  \def\CJKglue{}\def\CJK@kern{}\def\CJK@CJK{}%
  漢\hskip 0pt plus 3pt 漢\hskip 0pt plus 3pt 漢\hskip 0pt plus 3pt
  漢\hskip 0pt plus 3pt 漢\hskip 0pt plus 3pt 漢\hskip 0pt plus 3pt
  漢\hskip 0pt plus 3pt 漢\hskip 0pt plus 3pt 漢
  \hbox to 0.5\zw{、\hss}\hskip 0.5\zw plus 0.25\zw minus 0.5\zw
  漢漢漢漢漢漢漢漢漢}

% ---- Panel D: 補正カーンとグルーの間に whatsit
\panel{D}{%
  \def\CJKglue{}\def\CJK@kern{}\def\CJK@CJK{}%
  漢\hskip 0pt plus 3pt 漢\hskip 0pt plus 3pt 漢\hskip 0pt plus 3pt
  漢\hskip 0pt plus 3pt 漢\hskip 0pt plus 3pt 漢\hskip 0pt plus 3pt
  漢\hskip 0pt plus 3pt 漢\hskip 0pt plus 3pt 漢
  、\kern -0.5\zw \pdfsavepos\hskip 0.5\zw plus 0.25\zw minus 0.5\zw
  漢漢漢漢漢漢漢漢漢}

% ---- Panel E: 禁則ペナルティは阻害しない
\panel{E}{%
  \def\CJKglue{}\def\CJK@kern{}\def\CJK@CJK{}%
  漢\hskip 0pt plus 3pt 漢\hskip 0pt plus 3pt 漢\hskip 0pt plus 3pt
  漢\hskip 0pt plus 3pt 漢\hskip 0pt plus 3pt 漢\hskip 0pt plus 3pt
  漢\hskip 0pt plus 3pt 漢\hskip 0pt plus 3pt 漢
  \penalty 10000 、\kern -0.5\zw \hskip 0.5\zw plus 0.25\zw minus 0.5\zw
  漢漢漢漢漢漢漢漢漢}

\end{document}
